\documentclass[11pt]{article}
\usepackage[margin=1in]{geometry}
\usepackage{amsmath}
\usepackage{booktabs}
\usepackage{tabularx}
\usepackage{pgfplots}
\pgfplotsset{compat=1.18}
\usepackage{listings}
\usepackage{xcolor}
\usepackage{xurl}
\usepackage[hidelinks]{hyperref}

\lstset{
  basicstyle=\ttfamily\small,
  frame=single,
  columns=fullflexible,
  keepspaces=true,
  breaklines=true
}

\newcommand{\lecdate}{November 22, 2026}


\newcommand{\handout}[5]{
  \noindent
  \begin{center}
  \framebox{
    \vbox{
      \hbox to 5.78in { {\bf CSC 38000 - Introduction to Computer Security} \hfill #2 }
      \vspace{4mm}
      \hbox to 5.78in { {\Large \hfill #5  \hfill} }
      \vspace{1mm}
      \hbox to 5.78in { \hfill \lecdate \hfill }
      \vspace{3mm}
      \hbox to 5.78in { {\em #3} \hfill }
      \vspace{1mm}
      \hbox to 5.78in { {\em #4} \hfill }
    }
  }
  \end{center}
  \vspace*{-6mm}
}

\newcommand{\lecture}[4]{\handout{#1}{#2}{#3}{#4}{#1}}

\newtheorem{theorem}{Theorem}
\newtheorem{corollary}[theorem]{Corollary}
\newtheorem{lemma}[theorem]{Lemma}
\newtheorem{observation}[theorem]{Observation}
\newtheorem{proposition}[theorem]{Proposition}
\newtheorem{definition}[theorem]{Definition}
\newtheorem{claim}[theorem]{Claim}
\newtheorem{fact}[theorem]{Fact}
\newtheorem{assumption}[theorem]{Assumption}

% 1-inch margins, from fullpage.sty by H.Partl, Version 2, Dec. 15, 1988.
\topmargin 0pt
\advance \topmargin by -\headheight
\advance \topmargin by -\headsep
\textheight 8.9in
\oddsidemargin 0pt
\evensidemargin \oddsidemargin
\marginparwidth 0.5in
\textwidth 6.5in

\parindent 0in
\parskip 1.5ex

\begin{document}

\lecture{Progress Report: TOTP Security}{Fall 2026}{Lecturer: Garrison Grogan}{Name: Zeba Fariha, Nirath Hussan, Rashmi Khadka, Jacob Li}

\section{Overview}

Our research question is: \textbf{How does widening the TOTP acceptance window
affect the balance between successful legitimate authentication and resistance
to guessing and replay attacks?}\\
For this milestone (November 10 -- 22) we planned to build and test the TOTP
code generator and verifier, run a first experiment comparing two acceptance
windows, and report early results.

\section{What We Built}

The project is written in C++17. Each group member built one part.

\begin{center}
\begin{tabularx}{\textwidth}{l l X}
\toprule
Part & Files & What it does \\
\midrule
1. Generator & \texttt{totp.h}, \texttt{totp.cpp} &
  Turns the shared secret and the current time into a 6-digit code, following
  RFC 6238~\cite{rfc6238}. \\
  & \texttt{test\_totp.cpp} &
  Checks the generator against the 6 official RFC 6238 test codes. All pass. \\
2. Verifier & \texttt{verifier.h}, \texttt{verifier.cpp} &
  Server-side check of a code. Supports window 0 (current 30 s step only) and
  window 1 (previous, current, and next step). Rejects badly formed codes,
  blocks a code from being used twice, and can lock the account after too many
  wrong codes. \\
3. Experiments & \texttt{experiments.cpp} &
  Pretends the phone's clock is wrong by different amounts and records whether
  the server still accepts the code. Output: \texttt{results.csv}. \\
4. Demo and build & \texttt{main.cpp}, \texttt{Makefile} &
  A menu-based demo to generate and check codes, change the window, set a
  clock offset, and reset after a lockout. The \texttt{Makefile} builds
  everything. \\
  & \texttt{generate\_summary.ps1} &
  Script that turns \texttt{results.csv} into
  \texttt{acceptance\_summary.csv}: the acceptance rate for each window and
  clock offset. \\
\bottomrule
\end{tabularx}
\end{center}

\paragraph{How the code is made.} The generator divides the time by 30 to get
a step number, writes it as 8 bytes, computes HMAC-SHA-1 of it with the secret,
and cuts the 20-byte result down to 6 digits (dynamic truncation, RFC
4226~\cite{rfc4226}).

\paragraph{Our own HMAC.} Following the instructor's comment that core
cryptographic parts should not come from a library, we wrote HMAC ourselves
using the RFC 2104 formula~\cite{rfc2104}:
\[
  \text{HMAC}(K, M) = \text{SHA1}\big((K \oplus \text{opad}) \,\|\,
  \text{SHA1}((K \oplus \text{ipad}) \,\|\, M)\big).
\]
Only SHA-1 itself comes from the Crypto++ library. The RFC test codes still
pass, which shows our HMAC gives the correct output.

\section{How to Compile and Run}

\paragraph{Requirements.} A C++17 compiler and the Crypto++ library.
\begin{lstlisting}
# macOS
brew install cryptopp
# Ubuntu / Linux
sudo apt install libcrypto++-dev
\end{lstlisting}

\paragraph{Build and run.} From the project folder:
\begin{lstlisting}
make            # builds test_totp, experiments, and totp_demo
make test       # runs the RFC 6238 test codes and verifier checks
make results    # runs the experiment and rewrites results.csv
./totp_demo     # starts the demo (same command as in the COMMAND file)
make clean      # deletes the built programs
\end{lstlisting}
To rebuild the summary table from \texttt{results.csv}, run
\texttt{generate\_summary.ps1} in PowerShell (\texttt{pwsh} on macOS/Linux).

\paragraph{Using the demo.} The demo shows a menu:
\begin{lstlisting}
1. Generate TOTP code
2. Verify TOTP code
3. Change acceptance window
4. Set device clock offset
5. Reset settings
6. Exit
\end{lstlisting}
Example run with the phone clock 20 seconds behind the server:
\begin{lstlisting}
Option 4 -> offset -20     Option 1 -> Generated TOTP code: 807812
Option 3 -> window 0       Option 2 -> 807812 -> REJECTED (no match)
Option 3 -> window 1       Option 2 -> 807812 -> ACCEPTED (matched offset -1)
                           Option 2 -> 807812 -> REJECTED (reused)
\end{lstlisting}
The same clock error fails with window 0 but works with window 1, and using the
same code a second time is blocked. If a non-number is typed, the demo asks
again instead of closing.

\section{Preliminary Results}

\paragraph{Experiment setup.} The server checks a code at a fixed time. The
phone makes its code using a clock that is off by $-90$ to $+90$ seconds (in
5-second steps). We ran this for window 0 and window 1, and at 7 different
points inside the 30-second step (0, 5, 10, 15, 20, 25, 29 seconds), because
whether a code passes also depends on \emph{when} in the step the user logs
in. This gives 518 trials.

\paragraph{What this means.}
\begin{itemize}
  \item \textbf{Window 0 is strict.} A real user only gets in if their phone is
        in the same 30-second step as the server. Even a 5-second clock error
        fails in some trials, depending on timing.
  \item \textbf{Window 1 is forgiving.} Any clock error up to 30 seconds always
        works, and up to about 59 seconds sometimes works. Normal phone clocks
        are well inside this range.
  \item \textbf{The cost.} Window 1 accepts 3 codes at once instead of 1, so a
        random guess is 3 times more likely to work. How much this matters
        depends on attempt limits and reuse blocking, which we test next.
\end{itemize}


\begin{thebibliography}{9}
\bibitem{rfc6238} D. M'Raihi, S. Machani, M. Pei, and J. Rydell.
  \textit{TOTP: Time-Based One-Time Password Algorithm}. RFC 6238, IETF, May
  2011. \url{https://www.rfc-editor.org/info/rfc6238}
\bibitem{rfc4226} D. M'Raihi, M. Bellare, F. Hoornaert, D. Naccache, and
  O. Ranen. \textit{HOTP: An HMAC-Based One-Time Password Algorithm}. RFC 4226,
  IETF, December 2005. \url{https://www.rfc-editor.org/info/rfc4226}
\bibitem{rfc2104} H. Krawczyk, M. Bellare, and R. Canetti. \textit{HMAC:
  Keyed-Hashing for Message Authentication}. RFC 2104, IETF, February 1997.
  \url{https://www.rfc-editor.org/info/rfc2104}
\bibitem{nist} National Institute of Standards and Technology. \textit{Digital
  Identity Guidelines: Authentication and Authenticator Management}. NIST
  Special Publication 800-63B-4, 2025. \url{https://pages.nist.gov/800-63-4/sp800-63b/}
\end{thebibliography}

\end{document}
